\chapter{Related Work}

The gap this thesis addresses does not exist in isolation. It sits at the end
of a line of work on PDF processing tools, modern conversion systems, and
LLM-based document understanding. Each of these areas has pushed the field
forward, and each has a limitation that the next tried to fix. This chapter
traces that line and identifies where it stops.

\section{Traditional PDF Processing Tools}

Traditional PDF processing tools fall into two categories: text-layer parsers
that read content directly from the embedded glyph stream of a native PDF, and
OCR-based tools that recover text by applying image recognition to rendered
page images. Both are well-established, but they were built to extract text,
not to reconstruct the semantic structure of a document. That is the core
limitation both categories share.

\subsection{Text-based Parsers}

Text-layer parsers like PyMuPDF, pdfplumber, and pypdfium2 read character
objects and their coordinates from a PDF's internal object stream, then
reassemble them into reading order using heuristics. Each library works at a
different level of detail. PyMuPDF gives direct access to character bounding
boxes, font metadata, and color attributes, which makes it fast and useful as
a preprocessing step. pdfplumber is built on PDFMiner and exposes bounding-box
tables for words and characters, which helps with positional grouping for
tabular data. pypdfium2 wraps the PDFium rendering engine (the same one
used in Google Chrome) and produces high-fidelity rendering, but does
little structural post-processing on its own.

Meuschke et al.\ \autocite{meuschke2023} benchmarked ten tools on 500,000 pages
of academic documents and measured F1 scores across element types. The results
are specific. PyMuPDF scored F1=0.51 on paragraph extraction because it assigns
sections, captions, lists, and equations all to the same paragraph label, producing
noisy output. Table extraction was poor across every tool: the best result was
F1=0.47 (Adobe Extract), while dedicated table parsers like Camelot and Tabula
reached only 0.28--0.30, partly because they misidentify two-column layouts as
tables. List extraction scored zero for every tool tested. Equation extraction
was near zero. These are not edge cases --- they are common elements in academic
and technical documents. Ingemarsson and Persson \autocite{ingemarsson2024}
identified a further hard limit: on pages where content exists only as rasterised
images, text-layer parsers recover nothing. For mathematical expressions, Gao
et al.\ \autocite{gao2011} show that sequential character extraction loses the
spatial relationships that make formulas meaningful. The output is
character-correct, but the structure is gone.

\subsection{OCR Solutions}

OCR tools work differently: instead of reading a text layer, they apply image
recognition to rendered page images. Smith \autocite{smith2007} describes the
architecture of Tesseract, the most widely used open-source OCR engine. It
runs a two-pass pipeline: the first pass uses adaptive thresholding and
connected-component analysis to isolate character candidates, and the second
pass applies a word-level language model to resolve ambiguous boundaries.
Tesseract 4.0 added a Long Short-Term Memory (LSTM)-based recognition model that improved accuracy on
degraded and stylised text while keeping the same two-stage structure.
Subramani et al.\ \autocite{subramani2020} reviewed deep learning approaches to
document understanding and placed classical engines like Tesseract in that
broader context. Their conclusion is that single-pass OCR works well on clean,
single-font documents but breaks down on complex layouts, mixed-script content,
and low-resolution scans. This is what pushed the field toward layout-aware,
multi-stage pipelines.

Zhu et al.\ \autocite{zhu2022} developed DocBed, a multi-stage OCR pipeline built
for documents with complex layouts. Rather than applying recognition to the
full page at once, DocBed first runs a region segmentation model to separate
text blocks, tables, and figures, then sends each region to a dedicated
recogniser. This staged approach recovered significantly more structure than
global OCR on layout-complex documents. But it also introduced a known failure
mode: if the segmentation step makes a wrong decision, it passes incorrect
bounding boxes to every recogniser that follows, and the resulting errors are
hard to trace back to their source. This is not a problem unique to DocBed.
It applies to any modular pipeline architecture \autocite{duan2025}.

\subsection{Shared Limitations}

Both text-layer parsers and OCR tools share a gap that Darvishy
\autocite{darvishy2018} identified: no traditional tool can interpret or describe
mathematical formulas or scientific graphs in PDFs. Ouyang et al.\
\autocite{omnidocbench2025} show this extends even to specialised formula tools:
exact-match rates are near zero despite high character-level scores, because
structural grouping is lost even when individual symbols are correct. Figures
fare no better: every tool in their benchmark omitted figure content from
structured output entirely. Multi-column layouts are a third shared weakness:
reading-order error roughly doubles on complex layouts compared to single-column
pages across every tool category tested \autocite{omnidocbench2025}. For the
academic and scientific documents this thesis targets, formulas, figures, and
multi-column layouts are not edge cases. The evidence
from Meuschke et al.\ \autocite{meuschke2023}, Ingemarsson and Persson
\autocite{ingemarsson2024}, and Subramani et al.\ \autocite{subramani2020} all point
to the same conclusion: traditional tools work well within a narrow range of
document types, but they are not sufficient for the diverse, structure-rich
documents this thesis works with. This is what motivated the move toward
learning-based and multimodal approaches.

\section{Modern Document Conversion Tools}

Modern conversion tools go beyond simple text extraction. Rather than reading the
raw text layer, they first analyse the visual structure of each page. This approach
is sometimes called pipeline-based document parsing: the page is decomposed into
layout blocks, each classified as text, table, figure, or heading, and the blocks
are then processed in reading order \autocite{whitington2012}. Scanned pages are detected
automatically and routed to an OCR module, while digitally-born pages go through a
text extraction path. Subramani et al.\ \autocite{subramani2020} survey these methods,
distinguishing between modular pipeline systems and unified Vision-Language Model
(VLM) approaches. Robust handling of complex layouts remains an open challenge
for both.

Commercial tools such as LlamaParse, Marker, and Mistral OCR follow this general
architecture but add API layers and post-processing for downstream use cases.
Ouyang et al.\ \autocite{omnidocbench2025} evaluated pipeline-based and end-to-end
multimodal parsers across diverse real-world document layouts in OmniDocBench.
No single parser performed consistently well across all document types. Gaps were
largest for charts and degraded scanned pages. These results highlight how
difficult it is to build one general-purpose tool for structurally diverse
documents.

Blecher et al.\ \autocite{blecher2023} introduced Nougat, a Visual Transformer
trained on paired PDF-markup data that converts full academic document pages
directly to Markdown. Nougat removes the need for a separate layout analysis
or OCR stage by generating structured output end to end. It handles headings,
equations, tables, and citations from a single model pass. Subramani et al.\
\autocite{subramani2020} position this shift in a broader context: it represents
the move from pipeline-based recognition to unified visual models that reason
over the full page at once.

For long, multimodal academic documents, Xie et al.\ \autocite{wukong2024} introduced
PDF-WuKong, a multimodal large language model for question-answering over long
PDFs. The system uses a sparse sampler that operates on both text and image
representations, selecting only the paragraphs and diagrams most relevant to a
given query. Evaluated on the PaperPDF dataset of 1.1 million QA pairs from
English and Chinese academic papers, PDF-WuKong outperformed proprietary products
by an average of 8.6\% on F1. This shows that efficient joint modelling of text
and images is achievable for long documents, though the system targets retrieval
and question-answering rather than full document conversion.

Even so, key limitations remain. Meuschke et al.\ \autocite{meuschke2023} found
that even modern parsing tools struggle with scientific and patent documents
containing dense formulas and non-standard layouts.
Darvishy \autocite{darvishy2018} notes that no existing tool provides mechanisms to
interpret and describe mathematical formulas or scientific graphs in PDFs. This gap
is one of the central motivations for the image-based extraction strategy in this
thesis, where rendered page images are passed to a multimodal LLM that can reason
about text and visual content together.

\section{LLM-based Document Understanding}

Before large language models, document understanding depended on chains of
specialised tools. Subramani et al.\ \autocite{subramani2020} showed that even
deep-learning OCR broke down on complex layouts, mixed scripts, and
low-resolution pages. Vision-capable language models changed this: they read a
rendered page image and produce structured output token by token through
autoregressive generation, with no hand-coded layout rules in between.

Nougat \autocite{blecher2023}, described in the previous section, is the clearest
example of this shift: it infers headings, equations, tables, and citations in
a single model pass with no separate layout step. Meuschke et al.\
\autocite{meuschke2023} confirmed the pattern: structured-output tools
outperformed all rule-based parsers on scientific documents. A single
multimodal model can replace the multi-stage pipeline, at least for academic
documents.

OpenDataLab \autocite{mineru_software} released MinerU, an open-source dual-engine
system that automatically detects whether each page is scanned or digitally born
and routes it to a Vision-Language Model or a classical OCR path. It supports
109 languages and outputs structured Markdown or JSON (JavaScript Object Notation). Duan \autocite{duan2025}
pushes the argument further:
end-to-end decoder transformers outperform pipeline-based systems on accuracy
because they never pass an incorrect intermediate result down a fixed processing
chain.

Ding et al.\ \autocite{ding2026} survey Visually Rich Document Content Understanding
(VRD-CU), covering documents that combine text, layout, and visual elements.
Models pretrained across all three modalities consistently outperform text-only
approaches. Instruction-tuned LLM architectures now achieve cross-modal reasoning
that earlier models could not. Persistent gaps remain: word-level representations
miss document-level semantic structure, and multi-page documents with mixed
content are still an open challenge.

No current model covers all content types reliably. This gap motivates the
image-based strategy in this thesis, where rendered page images are passed to a
multimodal LLM and a postprocessing pipeline corrects structural failures in the
output.

Table~\ref{tab:tool_comparison} compares the tools discussed in this chapter
on the dimensions most relevant to structured document conversion.
\textit{Partial} in the Formulas column means formula pages are routed to a
vision model but the LaTeX output is not validated; \textit{Partial} in the
Routing column means the tool routes between two fixed paths rather than by
content type.

\begin{table}[ht]
\centering
\caption{Comparison of PDF processing tools across dimensions relevant to
structured document conversion}
\label{tab:tool_comparison}
\small\setlength{\tabcolsep}{4pt}
\renewcommand{\arraystretch}{1.3}
\begin{tabularx}{\textwidth}{lXXXXX}
\toprule
\textbf{Tool} & \textbf{Input} & \textbf{Formulas} & \textbf{Tables} & \textbf{Per-page routing} & \textbf{Output} \\
\midrule
PyMuPDF    & native  & No      & Partial & No              & text     \\
Tesseract  & scanned & No      & No      & No              & text     \\
Nougat     & native  & Yes     & Yes     & No              & Markdown \\
MinerU     & both    & Yes     & Yes     & Partial         & Markdown \\
\textbf{This work} & native & Partial & Yes & \textbf{Yes} & Markdown \\
\bottomrule
\end{tabularx}
\end{table}

\section{Evaluation Methods for Document Conversion}

Evaluating document conversion is harder than evaluating plain text extraction.
The output combines headings, paragraphs, tables, formulas, and figures, and each
element type can fail in a different way. Early work used string-similarity
metrics: BLEU score and edit distance (Levenshtein similarity) both compare
extracted text against a reference string, character by character
\autocite{papineni2002}. Meuschke et al.\ \autocite{meuschke2023} added local
alignment as a third metric, which allows gaps between matched segments and
catches word-order errors, such as merged columns or reversed reading order, that
BLEU misses.

All three metrics treat the full output as a flat sequence. A table with one
column shifted receives a low Levenshtein score even when nearly all the content
is correct. Tree Edit Distance based Similarity (TEDS) addresses this for tables:
it compares the predicted tree structure against a reference and gives partial
credit when individual cells are right even if the overall structure differs.
Research in document parsing evaluation has shown that no single metric covers
all content types, with layout analysis using mean Average Precision (mAP), text
extraction using edit metrics, and table evaluation using TEDS
\autocite{zhong2020, pfitzmann2022}.

Aggregate text-similarity scores hide capability gaps. A parser may score well
under BLEU or edit distance but have a misaligned table header or a dropped
superscript that breaks any downstream process. Ouyang et al.\
\autocite{omnidocbench2025} benchmarked diverse parsers across multiple content
dimensions and found that no system performed well on all types, confirming
that a single aggregate score cannot capture document parsing quality.

For conversion to Markdown specifically, N-gram metrics carry an additional
limitation: they ignore syntax. A heading marked with \texttt{\#} and the same
text written as bold have near-identical character content but different semantic
meaning in the rendered document. This is why the evaluation in Chapter~4 scores
structural, content, and formatting dimensions separately rather than collapsing
them into one number.
